\section{Perbandingan Efisiensi dan Studi Kasus}
\label{sec:efisiensi-studi-kasus}

Setelah seluruh mode pengalamatan dikenal, pertanyaan berikutnya
adalah mode mana yang sebaiknya dipakai. Jawabannya bergantung pada
jenis data yang diakses dan biaya yang bersedia dibayar dalam bentuk
ukuran kode serta waktu eksekusi.

\subsection{Ringkasan Tujuh Mode Pengalamatan}

Tabel~\ref{tab:tujuh-mode} merangkum ketujuh mode pengalamatan 8086.
Jumlah byte pada kolom terakhir dihitung untuk instruksi
\reg{AX} sebagai tujuan, dengan perpindahan 8 bit bila ada.

\begin{table}[H]
    \centering
    \caption{Tujuh mode pengalamatan 8086, segmen \eng{default}, dan
    jumlah byte instruksi contohnya.}
    \label{tab:tujuh-mode}
    \footnotesize
    \begin{tabular}{@{}p{1.8cm}p{2.8cm}p{3.6cm}p{2.1cm}p{1.5cm}@{}}
        \toprule
        \textbf{Mode} & \textbf{Bentuk penulisan} &
        \textbf{Contoh} & \textbf{Segmen default} &
        \textbf{Jumlah byte} \\
        \midrule
        Immediate & nilai konstan &
        \lstasm|MOV AX, 1234h| & tidak ada & 3 \\
        \addlinespace
        Register & register &
        \lstasm|MOV AX, BX| & tidak ada & 2 \\
        \addlinespace
        Direct & \code{[disp16]} atau nama variabel &
        \lstasm|MOV AX, [1234h]| & \reg{DS} & 3 \\
        \addlinespace
        Register indirect & \code{[BX]}, \code{[SI]}, \code{[DI]},
        \code{[BP]} & \lstasm|MOV AX, [BX]| &
        \reg{DS} (\reg{BP}: \reg{SS}) & 2 \\
        \addlinespace
        Indexed & \code{[SI+disp]}, \code{[DI+disp]} &
        \lstasm|MOV AX, [SI+8]| & \reg{DS} & 3 \\
        \addlinespace
        Base-relative & \code{[BX+disp]}, \code{[BP+disp]} &
        \lstasm|MOV AX, [BX+8]| &
        \reg{DS} (\reg{BP}: \reg{SS}) & 3 \\
        \addlinespace
        Based-indexed & \code{[BX+SI+disp]} dan tiga kombinasi lain &
        \lstasm|MOV AX, [BX+SI+8]| &
        \reg{DS} (\reg{BP}: \reg{SS}) & 3 \\
        \bottomrule
    \end{tabular}
\end{table}

\begin{infobox}
Jumlah byte tersebut adalah opcode ditambah byte \eng{ModR/M}
ditambah perpindahan. Bila perpindahannya 16 bit, mode indexed,
base-relative, dan based-indexed bertambah satu byte; bila
\eng{segment override prefix} dipakai, satu byte lagi ditambahkan.
Susunan byte ini diuraikan pada Bab~\ref{chap:format-instruksi}.
\end{infobox}

\subsection{Jumlah Byte dan Siklus Relatif}

Semakin banyak komponen alamat yang harus dijumlahkan, semakin lama
alamat efektif terbentuk. Tabel~\ref{tab:efisiensi-mode}
membandingkan ukuran kode dan waktu eksekusi setiap mode. Angka
siklus diambil dari tabel waktu Intel 8086 \cite{intel1979} dan
bersifat perkiraan karena eksekusi sesungguhnya juga dipengaruhi
antrean instruksi (\eng{instruction queue}) \cite{uva2024}.

\begin{table}[H]
    \centering
    \caption{Perbandingan efisiensi mode pengalamatan pada 8086.}
    \label{tab:efisiensi-mode}
    \footnotesize
    \begin{tabular}{@{}p{2.7cm}p{3.5cm}cp{1.6cm}p{4.4cm}@{}}
        \toprule
        \textbf{Mode} & \textbf{Contoh} & \textbf{Byte} &
        \textbf{Siklus} & \textbf{Paling cocok untuk} \\
        \midrule
        Register & \lstasm|MOV AX, BX| & 2 & 2 &
        data kerja yang paling sering \\
        \addlinespace
        Immediate & \lstasm|MOV AX, 1234h| & 3 & 4 &
        nilai awal dan tetapan \\
        \addlinespace
        Direct & \lstasm|MOV AX, [1234h]| & 3 & 10 &
        variabel tunggal ber-alamat tetap \\
        \addlinespace
        Register indirect & \lstasm|MOV AX, [BX]| & 2 & 13 &
        penelusuran \eng{pointer} \\
        \addlinespace
        Indexed & \lstasm|MOV AX, [SI+8]| & 3 & 17 &
        elemen array berindeks \\
        \addlinespace
        Base-relative & \lstasm|MOV AX, [BX+8]| & 3 & 17 &
        medan \eng{record} dan \code{[BP+disp]} \\
        \addlinespace
        Based-indexed & \lstasm|MOV AX, [BX+SI+8]| & 3 & 19 &
        matriks dan array di dalam \eng{record} \\
        \bottomrule
    \end{tabular}
\end{table}

Pertambahan kemampuan selalu dibayar dengan pertambahan waktu: mode
register selesai dalam dua siklus, sedangkan based-indexed hampir
sepuluh kali lipatnya. Karena itu, variabel yang paling sering dipakai
sebaiknya ditempatkan di register, dan mode berpenunjuk hanya dipakai
ketika alamat memang harus dihitung.

\subsection{Satu Operasi dalam Enam Mode Berbeda}

Kode~\ref{lst:satu-operasi-enam-mode} menuliskan satu operasi yang
sama dalam tujuh mode: menyalin elemen berindeks 2 array
\code{data1}, yaitu nilai \hex{1E}, ke register \reg{AX}.

\begin{lstlisting}[
    language=MASM,
    caption={Satu operasi yang sama dituliskan dalam tujuh mode
    pengalamatan berbeda.},
    label={lst:satu-operasi-enam-mode}
]
.DATA
    data1   DW  10, 20, 30, 40, 50   ; array, elemen ke-2 = 30 (1Eh)

.CODE
    ; 1. immediate        : nilainya sudah diketahui saat menulis kode
    MOV AX, 30                       ; AX = 1Eh

    ; 2. register         : nilainya sudah ada di register lain
    MOV BX, 30
    MOV AX, BX                       ; AX = 1Eh

    ; 3. direct           : alamat ditulis apa adanya
    MOV AX, data1+4                  ; AX = DS:(data1 + 4)

    ; 4. register indirect: offset disimpan di register penunjuk
    MOV BX, OFFSET data1
    ADD BX, 4
    MOV AX, [BX]                     ; AX = DS:BX

    ; 5. indexed          : [SI + perpindahan]
    MOV SI, 4
    MOV AX, data1[SI]                ; AX = DS:(data1 + SI)

    ; 6. base-relative    : [BX + perpindahan]
    MOV BX, OFFSET data1
    MOV AX, [BX+4]                   ; AX = DS:(BX + 4)

    ; 7. based-indexed    : [BX + SI + perpindahan]
    MOV BX, OFFSET data1
    MOV SI, 4
    MOV AX, [BX+SI]                  ; AX = DS:(BX + SI)
\end{lstlisting}

Ketujuh potongan itu menghasilkan nilai yang sama pada \reg{AX},
tetapi berbeda panjang kode dan jumlah siklusnya. Cara 4 sampai 7
lebih menguntungkan bila indeks berubah saat program berjalan, karena
hanya isi register yang perlu diubah.

\subsection{Studi Kasus: Elemen Array Satu Dimensi}

Array satu dimensi disimpan berurutan. Untuk array \code{word} 8086,
elemen ke-$i$ berada pada offset awal array ditambah $2i$, sehingga
seluruh array dapat ditelusuri dengan satu instruksi tetap.

\begin{lstlisting}[language=MASM, numbers=none]
    MOV BX, OFFSET data1     ; BX = basis, awal array
    XOR SI, SI               ; SI = 0, indeks elemen pertama
    MOV CX, 5                ; banyaknya elemen
ULANG:
    MOV AX, [BX+SI]          ; AX <- data1[i]
    ; ... olah nilai AX di sini ...
    ADD SI, 2                ; maju ke elemen berikutnya (2i + 2)
    LOOP ULANG               ; ulangi sampai CX = 0
\end{lstlisting}

Bila elemen array berukuran satu byte, langkahnya cukup
\lstasm|INC SI|; bila empat byte, digunakan \lstasm|ADD SI, 4|. Pola
yang berubah hanyalah besarnya langkah, bukan bentuk pengalamatannya.
Pola inilah yang dipakai pada \file{mode-array.asm} di
Bagian~\ref{sec:contoh-bab8}.

\subsection{Studi Kasus: Medan Record dengan Base-Relative}

\eng{Record} menyatukan beberapa medan bertipe berbeda dalam satu
blok memori. Karena panjang medannya tetap, offsetnya pun tetap:
register basis menunjuk awal \eng{record}, perpindahan memilih
medannya.

\begin{lstlisting}[language=MASM, numbers=none]
.DATA
    mhs1    DB  '2210'       ; medan NIM   : offset 0 sampai 3
    nilai1  DW  85           ; medan nilai : offset 4 sampai 5
.CODE
    MOV BX, OFFSET mhs1      ; basis = awal record
    MOV AX, [BX+4]           ; AX <- medan nilai (85)
\end{lstlisting}

Bentuk yang sama dengan register \reg{BP} dipakai untuk membaca
parameter prosedur karena \code{[BP+disp]} menunjuk segmen stack
secara bawaan (Bagian~\ref{sec:contoh-bab8} dan
Bab~\ref{chap:prosedur}).

\begin{studikasusbox}
Sebuah program menyimpan nilai mahasiswa dalam array \eng{record}
enam byte: empat byte NIM dan dua byte nilai. Alamat medan nilai
mahasiswa ke-$i$ dibentuk dengan \lstasm|[BX+SI+4]|, dengan \reg{BX}
menunjuk awal array dan \reg{SI} bernilai $6i$. Penulisan ini hanya
dua instruksi lebih panjang daripada \lstasm|MOV AX, [BX]|, tetapi
mampu menelusuri seluruh tabel tanpa mengubah instruksinya.
Perhatikan bahwa \reg{BX} menunjuk \reg{DS}, sehingga array nilai
harus berada pada segmen data, bukan segmen stack.
\end{studikasusbox}
